\chapter{Background}
\label{ch:background}

This chapter builds every concept the thesis depends on, in dependency
order: the GPU execution model (\S\ref{sec:bg-gpu}), the memory hierarchy and
its address spaces (\S\ref{sec:bg-mem}), why data movement dominates cost
(\S\ref{sec:bg-wall}), visual SLAM and cuVSLAM (\S\ref{sec:bg-slam}), the
candidate hardware substrates (\S\ref{sec:bg-substrates}), and the
measurement instruments (\S\ref{sec:bg-tools}).

\section{The GPU execution model}
\label{sec:bg-gpu}

A CPU optimizes the latency of one instruction stream; a GPU optimizes the
throughput of tens of thousands of them. NVIDIA GPUs are programmed through
CUDA: the programmer writes a \emph{kernel}---one function---and launches it
over a grid of thousands to millions of \emph{threads}, each processing a
different data element. Threads are grouped in fixed bundles of 32 called
\emph{warps}; the 32 \emph{lanes} of a warp execute the same instruction in
lockstep on different data. Warps are grouped into thread blocks, and blocks
are scheduled onto \emph{streaming multiprocessors} (SMs)---the GPU's core
unit. The primary instrument of this thesis, an NVIDIA RTX~2000 Ada
generation GPU (compute capability \code{sm\_89}), has 22~SMs, a 25\,MB
shared L2 cache, and 16\,GB of GDDR6 DRAM.

Two warp-level behaviours matter throughout this thesis:

\textbf{Divergence.} If a branch disables some lanes of a warp, the disabled
lanes do no useful work. The counter \code{mean\_active\_lanes} measures
this; 32.0 means fully converged. Every cuVSLAM kernel measured in this
thesis runs at 32.0 active lanes---divergence is \emph{not} cuVSLAM's
problem, which immediately focuses attention on memory behaviour.

\textbf{Coalescing.} When a warp issues a load, the hardware merges the 32
lane addresses into fetches of 32-byte \emph{sectors} (four sectors form a
128-byte cache line). If the lanes address 32 consecutive words, a few
sectors satisfy the whole warp---a \emph{coalesced} access. If the lanes
scatter, the hardware may fetch up to 32 sectors and discard most of each
one---a \emph{scattered gather} that wastes up to 8$\times$ the useful
bandwidth. The metric \emph{sectors per request} quantifies this ($\approx$4
= perfectly coalesced for 4-byte accesses; $\rightarrow$32 = fully
scattered). Coalescing has no CPU analogue and is one of the three reasons
DAMOV cannot be applied to GPUs unchanged (Chapter~\ref{ch:method}).

\emph{Occupancy} is the fraction of an SM's maximum concurrent warps that are
resident. It is the GPU's latency-hiding mechanism: while one warp waits for
memory, the scheduler issues another. High occupancy can hide DRAM
\emph{latency} entirely; nothing hides a \emph{bandwidth} shortfall. This
asymmetry---latency hidable, bandwidth not---shapes the entire bottleneck
taxonomy of Chapter~\ref{ch:method}.

\section{The memory hierarchy and the five address spaces}
\label{sec:bg-mem}

GPU storage forms the usual pyramid: per-thread \emph{registers} (fastest,
smallest); per-SM on-chip storage of 100\,KB serving as L1 cache and
programmer-managed \emph{shared memory}; the chip-wide L2 cache (25\,MB); the
GPU's DRAM (16\,GB, reachable at a measured 205\,GB/s on the locked
instrument, roughly 100$\times$ the latency of a register); host RAM across
the PCIe bus; and persistent storage. Each level downward costs roughly an
order of magnitude more energy per byte.

Independent of this physical pyramid, every memory \emph{instruction}
addresses one of several logical \emph{spaces}, and confusing them invalidates
analyses (as Chapter~\ref{ch:locality} demonstrates on this very project):

\begin{description}
\item[Global] the program's real data in DRAM---everything allocated with
\code{cudaMalloc}. Instructions: \code{LDG}/\code{STG}. All DRAM-traffic and
locality claims in this thesis are about global space.
\item[Shared] the on-chip per-SM scratchpad (\code{LDS}/\code{STS}). Used
for convolution tiles, sort staging, reduction trees. \emph{Not DRAM}:
traffic here says nothing about the memory wall.
\item[Local] a per-thread region that is physically \emph{in DRAM} but
logically compiler scratch: when a kernel needs more temporaries than its
register budget, the compiler \emph{spills} them here
(\code{LDL}/\code{STL}). Local addresses are interleaved across threads, so
spill traffic is \emph{perfectly coalesced by construction}---a fact that
produced this project's most instructive measurement artifact
(\S\ref{sec:loc-correction}).
\item[Constant, texture] small read-only parameter storage, and image reads
through dedicated texture units. Texture fetches do not appear in the
address-trace instrument used here; the implications are bounded in
\S\ref{sec:attr-residuals}.
\end{description}

\section{The memory wall}
\label{sec:bg-wall}

Since the mid-1990s, processor arithmetic throughput has grown far faster
than memory bandwidth or latency has improved---the \emph{memory
wall}~\cite{wulf1995memorywall}. On contemporary hardware, a double-precision
multiply costs picojoules while a DRAM fetch of its operands costs hundreds
of picojoules to nanojoules; crossing PCIe to host memory or storage costs
more still. For low-arithmetic-intensity codes---image pipelines are the
canonical case---the processor mostly waits. The roofline
model~\cite{williams2009roofline} makes this quantitative
(\S\ref{sec:met-roofline}): performance is bounded by
$\min(\text{peak compute},\ \text{AI}\times\text{peak bandwidth})$, where
arithmetic intensity AI is FLOPs per DRAM byte. A kernel under the sloped
part of that roof cannot be helped by more arithmetic units---only by moving
fewer bytes, or moving them a shorter distance. The latter option is the
premise of memory-centric architecture.

\section{Visual SLAM and cuVSLAM}
\label{sec:bg-slam}

Visual SLAM estimates a camera's trajectory while building a map of the
environment~\cite{cadena2016slam,durrant2006slam}. The canonical pipeline has
four stages:

\begin{enumerate}
\item \textbf{Front-end.} Each incoming frame is converted to grayscale;
an \emph{image pyramid} (the frame at successively halved resolutions) and
gradient images are built; distinctive corners are detected (Good Features
To Track, the Shi--Tomasi criterion) and tracked frame-to-frame with
Lucas--Kanade optical flow. This stage runs on every frame at camera rate.
\item \textbf{Visual odometry.} From tracked features, the camera's motion
is estimated by solving a nonlinear least-squares problem.
\item \textbf{Bundle adjustment (BA).} Poses and 3-D landmark positions over
a window are refined jointly. Numerically, each iteration builds and solves
a linear system (a Hessian reduced by the Schur complement to eliminate
landmark variables); this thesis names that structure
\code{ba\_linear\_system} and measures precisely which kernels stream it.
\item \textbf{SLAM layer.} Selected frames become \emph{keyframes}; their
feature descriptors enter a \emph{keyframe database}. When the camera
revisits a mapped place, a database scan detects the \emph{loop closure} and
the accumulated drift of the whole trajectory is corrected. Loop closure is
rare (it fires on revisits) but is what makes long-session maps consistent
---and its database is the only data structure that \emph{grows} with
session length, a fact of central architectural consequence
(\S\ref{sec:attr-static}).
\end{enumerate}

cuVSLAM~\cite{cuvslam2025} is NVIDIA's production implementation: GPU-resident
front-end and solvers, a host-side keyframe/landmark store backed by LMDB (a
memory-mapped embedded database), and support for stereo, RGB-D (color plus
depth), monocular, and inertial (IMU-fused) camera configurations, in
odometry-only or full-SLAM modes, with synchronous and asynchronous SLAM
variants. This thesis treats cuVSLAM strictly as a measured artifact: its
algorithms are never modified; instrumentation is additive and switchable,
and Chapter~\ref{ch:validity} proves it output-neutral.

Accuracy of a SLAM system is scored against \emph{ground truth} (the true
trajectory from motion capture, GPS/INS, or an external tracker) using two
standard metrics~\cite{sturm2012tum,zhang2018tutorial}: \emph{absolute
pose/trajectory error} (APE/ATE)---the RMS position difference after optimal
alignment (SE(3) for metric-scale systems; Sim(3) when monocular scale is
free~\cite{umeyama1991}), and \emph{relative pose/trajectory error}
(RPE/RTE)---drift over fixed-length segments, expressed as a percentage,
which is scale-independent and therefore comparable across
kilometre-scale (KITTI) and room-scale (EuRoC, TUM) trajectories.

\section{Candidate substrates for memory-centric execution}
\label{sec:bg-substrates}

The thesis evaluates each kernel's fit against five homes:

\textbf{GPU (status quo).} Best when there is real arithmetic per byte or
reuse the caches can capture.

\textbf{CPU/host.} Best for tiny, launch-overhead-dominated kernels whose GPU
occupancy is negligible.

\textbf{DRAM-PiM (near-bank).} Compute placed at DRAM banks exposes internal
bandwidth several times the external interface and removes the
interconnect-energy toll per byte. Commercially realized in UPMEM's
DPUs~\cite{gomezluna2021upmem}, Samsung's HBM-PIM (function-in-memory
DRAM)~\cite{kwon2021hbmpim}, and SK~hynix's GDDR6-AiM~\cite{lee2022aim}.
The natural tenant is \emph{streaming, bandwidth-bound} work: much data,
little arithmetic, no reuse.

\textbf{Scatter-capable PiM.} Irregular gather/scatter defeats both caches
and coalescing; near-memory engines that issue fine-grained random accesses
(in the lineage of Tesseract and PIM-enabled
instructions~\cite{ahn2015tesseract,ahn2015pei}) are the architectural answer
when the access pattern cannot be fixed in software.

\textbf{Near-sensor SRAM / ISP.} Per-frame pixel processing can execute
adjacent to the image sensor, consuming the stream before it is ever written
to DRAM. The measured 1.68\,MB-per-frame host-to-device upload
(\S\ref{sec:char-transfers}) is exactly the traffic such placement removes.

\textbf{In/near-storage processing.} When a growing database is scanned
episodically, the scan can execute inside the storage
device~\cite{do2013smartssd,gu2016biscuit}, returning matches rather than
streaming the store through host memory, PCIe, and DRAM.

\section{The measurement instruments}
\label{sec:bg-tools}

No single tool sees everything; the methodology composes four views.

\textbf{Nsight Systems (nsys)} records the timeline: every kernel launch,
its duration, every PCIe copy, with near-zero overhead. It provides the
screening signal (which kernels matter) and, through NVTX ranges, the
kernel-to-pipeline-stage mapping.

\textbf{Nsight Compute (ncu)} reads hardware performance counters per kernel
by \emph{replaying} each kernel multiple times. It is precise about
\emph{how much} (bytes, hit rates, stall cycles) but blind to \emph{which
addresses}. Replay makes it expensive, so captures are bounded to windows
and to a curated metric set of $\sim$15--30 counters; requesting the full
set is infeasible on small-memory GPUs.

\textbf{NVBit}~\cite{villa2019nvbit} rewrites GPU machine code at load time
to invoke a callback on every memory instruction, yielding the exact
per-warp address stream---ground truth beneath ncu's aggregates---at
100--1000$\times$ slowdown. The tool used here (\code{mem\_trace}, extended
with launch-window and kernel-name gating, and with an allocation-event
sidecar) makes gigabyte-scale traces feasible and, joined with the
allocation journal of Chapter~\ref{ch:attribution}, attributable.

\textbf{NVTX + TaggedAllocator} name things. cuVSLAM ships with NVTX stage
annotations that are compiled out by default at two independent gates; this
project's instrumentation patch enables them, making the kernel-to-stage
table a \emph{measured} artifact rather than a guess from kernel names. The
TaggedAllocator (Chapter~\ref{ch:attribution}) journals every GPU allocation
with its call stack, giving each buffer a durable data-structure name.

A standing epistemic ladder orders these instruments: \emph{counters say how
much; traces say which addresses; the allocation journal says which data
structure; NVTX says which pipeline stage.} Every claim in this thesis cites
the highest rung available, and the two places where rungs disagreed became
findings (\S\ref{sec:loc-correction}, \S\ref{sec:attr-spaces}).
